\Needspace{6\baselineskip}\Article{351}
Behalten wir also sämtliche Bezeichnungen des vorigen Artikels bei, so ist klar, dass sich jeder Koeffizient der Gleichung, deren Wurzeln die \(\beta\) Summen \((\gamma,\lambda),\allowbreak (\gamma,\lambda^{\prime}),\allowbreak (\gamma,\lambda^{\prime\prime})\) usw. sind, auf eine Form wie \[A+a(\beta\gamma,1)+a^{\prime}(\beta\gamma,g)\ldots+a^\varepsilon(\beta\gamma,g^{\alpha-1})\] zurückführen lässt und alle Zahlen \(A,\allowbreak a\) usw. ganze Zahlen sind. Die Gleichung, deren Wurzeln die \(\beta\) Perioden mit \(\gamma\) Gliedern sind, die in einer anderen Periode \((\beta\gamma,k\lambda)\) enthalten sind, geht aus jener hervor, wenn man überall in den Koeffizienten jede beliebige Periode \((\beta\gamma,\mu)\) durch \((\beta\gamma,k\mu)\) ersetzt. Ist also \(\alpha=1\), so werden sämtliche \(\beta\) Perioden mit \(\gamma\) Gliedern durch eine Gleichung vom Grad \(\beta\) bestimmt, deren einzelne Koeffizienten sich auf die Form \(A+a(\beta\gamma,1)\) bringen lassen \textit{und daher bekannte Größen sind}, da \((\beta\gamma,1)=(n-1,1)=-1\) gilt. Ist dagegen \(\alpha>1\), so sind die Koeffizienten der Gleichung, deren Wurzeln sämtliche Perioden mit \(\gamma\) Gliedern sind, die in einer gegebenen Periode mit \(\beta\gamma\) Gliedern enthalten sind, bekannte Größen, sobald die Zahlenwerte aller \(\alpha\) Perioden mit \(\beta\gamma\) Gliedern bekannt sind. — Übrigens lässt sich die Berechnung der Koeffizienten dieser Gleichungen oft bequemer ausführen, besonders wenn \(\beta\) nicht sehr klein ist, indem man zuerst die Potenzsummen der Wurzeln ermittelt und daraus nach dem Newtonschen Satz die Koeffizienten ableitet, ebenso wie oben in Artikel 349.

\textit{Beispiel I.} Gesucht ist für \(n=19\) die Gleichung, deren Wurzeln die Summen \((6,1),\allowbreak (6,2),\allowbreak (6,4)\) sind. Bezeichnen wir diese Wurzeln jeweils mit \(p,\allowbreak p^{\prime},\allowbreak p^{\prime\prime}\) und die gesuchte Gleichung mit \[x^3-Axx+Bx-C=0\] so folgt \[A=p+p^{\prime}+p^{\prime\prime},\quad B=pp^{\prime}+pp^{\prime\prime}+p^{\prime}p^{\prime\prime},\quad C=pp^{\prime}p^{\prime\prime}\] Daher ist \[A=(18,1)=-1\] Ferner gilt
\SourceLeaf{436}
\[\begin{gathered}
pp^{\prime}=p+2p^{\prime}+3p^{\prime\prime},\quad pp^{\prime\prime}=2p+3p^{\prime}+p^{\prime\prime},\\
p^{\prime}p^{\prime\prime}=3p+p^{\prime}+2p^{\prime\prime}
\end{gathered}\] und damit \[B=6(p+p^{\prime}+p^{\prime\prime})=6(18,1)=-6\] Schließlich erhält man \[\begin{aligned}
C=(p+2p^{\prime}+3p^{\prime\prime})p^{\prime\prime}&=3(6,0)+11(p+p^{\prime}+p^{\prime\prime})\\
&=18-11=7
\end{aligned}\] Daher lautet die gesuchte Gleichung \[x^3+xx-6x-7=0\] Nach der anderen Methode haben wir \[p+p^{\prime}+p^{\prime\prime}=-1\] \[\begin{gathered}
pp=6+2p+p^{\prime}+2p^{\prime\prime},\quad p^{\prime}p^{\prime}=6+2p^{\prime}+p^{\prime\prime}+2p,\\
p^{\prime\prime}p^{\prime\prime}=6+2p^{\prime\prime}+p+2p^{\prime}
\end{gathered}\] und somit \[pp+p^{\prime}p^{\prime}+p^{\prime\prime}p^{\prime\prime}=18+5(p+p^{\prime}+p^{\prime\prime})=13\] sowie ebenso \[p^3+p^{\prime3}+p^{\prime\prime3}=36+34(p+p^{\prime}+p^{\prime\prime})=2\] Daraus wird nach dem Newtonschen Satz dieselbe Gleichung wie zuvor abgeleitet.

II. Gesucht ist für \(n=19\) die Gleichung, deren Wurzeln die Summen \((2,1),\allowbreak (2,7),\allowbreak (2,8)\) sind. Bezeichnet man diese jeweils mit \(q,\allowbreak q^{\prime},\allowbreak q^{\prime\prime}\), so findet man \[\begin{gathered}
q+q^{\prime}+q^{\prime\prime}=(6,1),\\
qq^{\prime}+qq^{\prime\prime}+q^{\prime}q^{\prime\prime}=(6,1)+(6,4),\quad qq^{\prime}q^{\prime\prime}=2+(6,2)
\end{gathered}\] Mit den Bezeichnungen des vorigen Beispiels ist daher die gesuchte Gleichung \[x^3-pxx+(p+p^{\prime\prime})x-2-p^{\prime}=0\] Die Gleichung, deren Wurzeln die Summen \((2,2),\allowbreak (2,3),\allowbreak (2,5)\) sind, die in \((6,2)\) enthalten sind, geht aus der vorhergehenden hervor, wenn man für \(p,\allowbreak p^{\prime},\allowbreak p^{\prime\prime}\) jeweils \(p^{\prime},\allowbreak p^{\prime\prime},\allowbreak p\) einsetzt; führt man dieselbe Substitution nochmals aus, so entsteht die Gleichung, deren Wurzeln die Summen \((2,4),\allowbreak (2,6)\). \((2,9)\) sind, die in \((6,4)\) enthalten sind.

\Needspace{9\baselineskip}\begin{center}\textit{Auf den vorhergehenden Untersuchungen beruht die Lösung der Gleichung \(X=0\).}\end{center}


\Needspace{6\baselineskip}\Article{352}
Die vorstehenden Sätze mit den angefügten Folgerungen enthalten die Hauptpunkte der gesamten Theorie, und das Verfahren zur Bestimmung der Werte der Wurzeln \(\Omega\) kann nun kurz angegeben werden.
\SourceLeaf{437}

Zuerst wähle man eine Zahl \(g\), die für den Modul \(n\) eine primitive Wurzel ist, und ermittle die kleinsten Reste der Potenzen von \(g\) bis \(g^{n-2}\) nach dem Modul \(n\). Man zerlege \(n-1\) in Faktoren, und zwar in Primfaktoren, wenn man das Problem auf Gleichungen möglichst niedrigen Grades zurückführen will. Diese seien in völlig beliebiger Reihenfolge \(\alpha,\allowbreak \beta,\allowbreak \gamma\ldots\zeta\), und man setze \[\frac{n-1}{\alpha}=\beta\gamma\ldots\zeta=a,\qquad\frac{n-1}{\alpha\beta}=\gamma\ldots\zeta=b,\quad\text{usw.}\] Man teile sämtliche Wurzeln \(\Omega\) in \(\alpha\) Perioden mit \(a\) Gliedern ein; jede einzelne von diesen wiederum in \(\beta\) Perioden mit \(b\) Gliedern, jede einzelne von diesen erneut in \(\gamma\) Perioden usw. Nach dem vorigen Artikel suche man die Gleichung vom Grad \(\alpha\), \((A)\), deren Wurzeln jene \(\alpha\) Summen mit \(a\) Gliedern sind; ihre Werte werden also durch Lösung dieser Gleichung bekannt.

Hier entsteht aber eine Schwierigkeit, da ungewiss erscheint, welcher Wurzel der Gleichung \((A)\) die einzelnen Summen gleichzusetzen sind, nämlich welche Wurzel mit \((a,1)\), welche mit \((a,g)\) usw. bezeichnet werden soll. Dies lässt sich folgendermaßen beheben. Mit \((a,1)\) kann man eine beliebige Wurzel der Gleichung \((A)\) bezeichnen. Denn jede Wurzel dieser Gleichung ist eine Summe aus \(a\) Wurzeln von \(\Omega\), und es ist völlig beliebig, welche Wurzel aus \(\Omega\) mit \([1]\) bezeichnet wird. Daher dürfen wir offenbar annehmen, dass irgendeine der Wurzeln, aus denen eine beliebige gegebene Wurzel der Gleichung \((A)\) besteht, durch \([1]\) dargestellt wird; dann wird jene Wurzel der Gleichung \((A)\) zu \((a,1)\). Die Wurzel \([1]\) ist dadurch aber noch nicht vollständig bestimmt; es bleibt weiterhin völlig beliebig oder unbestimmt, welche der Wurzeln, die \((a,1)\) bilden, wir für \([1]\) nehmen wollen. Sobald aber \((a,1)\) bestimmt ist, lassen sich daraus auch alle übrigen Summen mit \(a\) Gliedern rational ableiten (Artikel 346). Daraus geht zugleich hervor, dass man durch Lösung dieser Gleichung nur eine einzige Wurzel zu ermitteln braucht. — Für diesen Zweck kann auch das folgende, weniger direkte Verfahren verwendet werden. Man wähle für \([1]\) eine bestimmte Wurzel, \textit{das heißt}, man setze \([1]=\cos\frac{kP}{n}+i\sin\frac{kP}{n}\), wobei man die ganze Zahl \(k\) beliebig, jedoch so wählt, dass sie durch \(n\) nicht teilbar ist. Dann bezeichnen auch \([2],\allowbreak [3]\) usw. bestimmte Wurzeln und daher auch die Summen \((a,1),\allowbreak (a,g)\) usw. bestimmte Größen. Berechnet man diese mit leichter Feder aus Sinustafeln, nämlich so genau, dass sich entscheiden lässt, welche größer und welche kleiner sind, so kann kein Zweifel bleiben, durch welche Bezeichnungen die einzelnen Wurzeln der Gleichung \((A)\) zu unterscheiden sind.

Sind auf diese Weise sämtliche \(\alpha\) Summen mit \(a\) Gliedern gefunden, so suche man nach dem vorigen Artikel die Gleichung \((B)\) vom Grad \(\beta\), deren Wurzeln die \(\beta\) Summen mit \(b\) Gliedern sind, die in \((a,1)\) enthalten sind; sämtliche Koeffizienten dieser Gleichung sind bekannte
\SourceLeaf{438}
Größen. Da noch beliebig ist, welche der \(a=\beta b\) Wurzeln unter \((a,1)\) mit \([1]\) bezeichnet wird, kann eine beliebige gegebene Wurzel der Gleichung \((B)\) durch \((b,1)\) ausgedrückt werden, weil man offenbar annehmen darf, dass eine der \(b\) Wurzeln, aus denen sie besteht, mit \([1]\) bezeichnet wird. Man suche also durch Lösung der Gleichung \((B)\) eine beliebige Wurzel, setze sie \(=(b,1)\) und leite daraus nach Artikel 346 alle übrigen Summen mit \(b\) Gliedern ab. Dadurch erhalten wir zugleich eine Probe für die Rechnung, da diejenigen Summen mit \(b\) Gliedern, die zu denselben Perioden mit \(a\) Gliedern gehören, stets bekannte Summen ergeben müssen. — In einigen Fällen kann es ebenso einfach sein, \(\alpha-1\) weitere Gleichungen vom Grad \(\beta\) zu ermitteln, deren Wurzeln jeweils die einzelnen \(\beta\) Summen mit \(b\) Gliedern sind, die in den übrigen Perioden mit \(a\) Gliedern, \((a,g),\allowbreak (a,gg)\) usw., enthalten sind, und \textit{alle} Wurzeln sowohl dieser Gleichungen als auch der Gleichung \(B\) durch Lösung zu bestimmen. Dann muss man aber ebenso wie oben mithilfe einer Sinustafel entscheiden, welchen Perioden mit \(b\) Gliedern die einzelnen so erhaltenen Wurzeln gleichzusetzen sind. Für diese Entscheidung können auch verschiedene andere Kunstgriffe verwendet werden, die hier nicht vollständig erläutert werden können; einen für den Fall \(\beta=2\), der besonders nützlich ist und sich durch Beispiele kürzer als durch Regeln erläutern lässt, wird man jedoch in den folgenden Beispielen kennenlernen.

Sind auf diese Weise die Werte aller \(\alpha\beta\) Summen mit \(b\) Gliedern gefunden, so können daraus auf ganz ähnliche Weise durch Gleichungen vom Grad \(\gamma\) sämtliche \(\alpha\beta\gamma\) Summen mit \(c\) Gliedern bestimmt werden. Man muss nämlich \textit{entweder eine} Gleichung vom Grad \(\gamma\), deren Wurzeln die \(\gamma\) Summen mit \(c\) Gliedern sind, die unter \((b,1)\) enthalten sind, nach Artikel 350 ermitteln, durch ihre Lösung eine beliebige Wurzel finden und \(=(c,1)\) setzen und schließlich daraus nach Artikel 346 alle übrigen gleichartigen Summen ableiten; \textit{oder} ebenso insgesamt \(\alpha\beta\) Gleichungen vom Grad \(\gamma\) aufstellen, deren Wurzeln jeweils die \(\gamma\) Summen mit \(c\) Gliedern in den einzelnen Perioden mit \(b\) Gliedern sind, die Werte sämtlicher Wurzeln aller dieser Gleichungen durch Lösung bestimmen und schließlich ihre Reihenfolge ebenso wie oben mithilfe einer Sinustafel oder für \(\gamma=2\) durch den unten in den Beispielen zu zeigenden Kunstgriff festlegen.

Fährt man auf diese Weise fort, so erhält man offenbar schließlich sämtliche \(\frac{n-1}{\zeta}\) Summen mit \(\zeta\) Gliedern. Stellt man daher nach Artikel 348 die Gleichung vom Grad \(\zeta\) auf, deren Wurzeln die \(\zeta\) Wurzeln aus \(\Omega\) sind, die in \((\zeta,1)\) enthalten sind, so sind ihre Koeffizienten sämtlich bekannte Größen. Findet man durch Lösung eine beliebige Wurzel dieser Gleichung, so darf man diese \(=[1]\) setzen, und alle übrigen Wurzeln \(\Omega\) werden durch deren Potenzen
\SourceLeaf{439}
erhalten. Wenn man will, können auch alle Wurzeln jener Gleichung durch Lösung ermittelt werden und außerdem durch Lösung von \(\frac{n-1}{\zeta}-1\) weiteren Gleichungen vom Grad \(\zeta\), die jeweils alle \(\zeta\) Wurzeln in den einzelnen übrigen Perioden mit \(\zeta\) Gliedern darstellen, alle übrigen Wurzeln \(\Omega\) gefunden werden.

Übrigens ist klar, dass, sobald die erste Gleichung \((A)\) gelöst ist oder die Werte aller \(\alpha\) Summen mit \(a\) Gliedern bekannt sind, nach Artikel 348 unmittelbar auch die Zerlegung der Funktion \(X\) in \(\alpha\) Faktoren vom Grad \(a\) gegeben ist. Ferner wird nach Lösung der Gleichung \((B)\) oder nachdem die Werte aller \(\alpha\beta\) Summen mit \(b\) Gliedern gefunden sind, jeder jener Faktoren wiederum in \(\beta\) oder \(X\) in \(\alpha\beta\) Faktoren vom Grad \(b\) zerlegt usw.


